% META: {"id": "TCOMPL-MF-CO-008", "chapitre": "TCOMPL-MODELES-FONCTION", "type_objet": "corrige", "exercice_ref": "TCOMPL-MF-EX-008", "capacites_codes": ["C4"], "sources_inspiration": [], "status": "generated"}
% BEGIN-VERIFY
% from sympy import symbols, diff, solve
% x = symbols('x', real=True)
% f = x ** 3 - 3 * x
% assert diff(f, x) == 3 * x ** 2 - 3
% assert diff(f, x, 2) == 6 * x
% assert solve(diff(f, x, 2), x) == [0]
% assert diff(f, x, 2).subs(x, -1) < 0
% assert diff(f, x, 2).subs(x, 1) > 0
% assert f.subs(x, 0) == 0
% END-VERIFY
\begin{corrige}{TCOMPL-MF-EX-008}
\textbf{1.} La courbe tourne sa concavité vers le bas sur $[-2\,;\,0]$ (elle
y est concave) et vers le haut sur $[0\,;\,2]$ (convexe).

\textbf{2.} Le changement de courbure a lieu en $x=0$ : c'est le point
d'inflexion, de coordonnées $(0\,;\,0)$.

\textbf{3.} $f'(x)=3x^2-3$ et $f''(x)=6x$, qui s'annule en changeant de signe
en $x=0$ : négatif avant (concavité), positif après (convexité). Le calcul
confirme la lecture graphique.
\end{corrige}
